
% BEGIN S2
 
\par\medskip\noindent\textbf{2 Preliminaries}\par
 
% BEGIN S2.p1
 
In this section we fix preliminaries and basic notation, following standard literature [2].\par
 
% END S2.p1
 
% BEGIN S2.SS0.SSS0.Px1
 
\par\medskip\noindent\textbf{Higher-order $\beta$-Matching in the Simply Typed $\lambda$-Calculus}\par
 
% BEGIN S2.SS0.SSS0.Px1.p1
 
The syntax of untyped $\lambda$-terms is given in the following Definition 1.\par
 
% END S2.SS0.SSS0.Px1.p1
 
% BEGIN Thmtheorem1
 
\par\medskip\noindent\textbf{Definition 1 ($\lambda$-Terms).}\par
 
% BEGIN Thmtheorem1.p1
 
$M,N::=x\mid M\,N\mid\lambda x.M$ where $a,\ldots,z$ range over term variables.\par
 
% END Thmtheorem1.p1
 
% END Thmtheorem1
 
% BEGIN S2.SS0.SSS0.Px1.p2
 
Substitution of the term variable $x$ in the term $M$ by the term $N$ is denoted $M[x:=N]$.As usual, term application associates to the left, and we may group consecutive $\lambda$-abstractions. We commonly refer to the term $\lambda x.x$ as $I$.\par
 
% END S2.SS0.SSS0.Px1.p2
 
% BEGIN Thmtheorem2
 
\par\medskip\noindent\textbf{Definition 2 ($\beta$-Reduction).}\par
 
% BEGIN Thmtheorem2.p1
 
The relation $\to_{\beta}$ on terms is the contextual closure of $(\lambda x.M)\,N\to_{\beta}M[x:=N]$.\par
 
% END Thmtheorem2.p1
 
% END Thmtheorem2
 
% BEGIN S2.SS0.SSS0.Px1.p3
 
The $\beta$-equivalence relation $=_{\beta}$ is the reflexive, transitive, symmetric closure of $\to_{\beta}$.\par
 
% END S2.SS0.SSS0.Px1.p3
 
% BEGIN S2.SS0.SSS0.Px1.p4
 
In the simply typed $\lambda$-calculus we may assign to a term $M$ a simple type $\tau$ in type environment $\Gamma$, written $\Gamma\vdash M:\tau$. Similarly to prior work [13], one ground atom $\iota$ in the construction of simple types suffices for the negative result in the present work. Definition 5 contains the rules (Var), ($\to$I), and ($\to$E) of the simple type system.\par
 
% END S2.SS0.SSS0.Px1.p4
  
% BEGIN Thmtheorem3
 
\par\medskip\noindent\textbf{Definition 3 (Simple Types with Ground Atom $\iota$).}\par
 
% BEGIN Thmtheorem3.p1
 
\par
    \[\sigma,\tau::=\iota\mid\sigma\to\tau\]    
% END Thmtheorem3.p1
 
% END Thmtheorem3
 
% BEGIN S2.SS0.SSS0.Px1.p5
 
The arrow type constructor $\to$ associates to the right.\par
 
% END S2.SS0.SSS0.Px1.p5
 
% BEGIN Thmtheorem4
 
\par\medskip\noindent\textbf{Definition 4 (Type Environments).}\par
 
% BEGIN Thmtheorem4.p1
    \[\Gamma::=\{x_{1}:\sigma_{1},\ldots,x_{n}:\sigma_{n}\}\text{ where }x_{i}\neq x%
_{j}\text{ for }i\neq j\]    
% END Thmtheorem4.p1
 
% END Thmtheorem4
 
% BEGIN Thmtheorem5
 
\par\medskip\noindent\textbf{Definition 5 (Simple Type System).}\par
 
% BEGIN Thmtheorem5.p1
 
\begin{center}\includegraphics[width=0.95\linewidth,height=0.45\textheight,keepaspectratio]{formula-01.pdf}\end{center}\par
 
% END Thmtheorem5.p1
 
% END Thmtheorem5
 
% BEGIN S2.SS0.SSS0.Px1.p6
 
The following Example 6 illustrates a type derivation in the simple type system.\par
 
% END S2.SS0.SSS0.Px1.p6
 
% BEGIN Thmtheorem6
 
\par\medskip\noindent\textbf{Example 6.}\par
 
% BEGIN Thmtheorem6.p1
 
\begin{center}\includegraphics[width=0.95\linewidth,height=0.45\textheight,keepaspectratio]{formula-02.pdf}\end{center}
 
% END Thmtheorem6.p1
 
% END Thmtheorem6
 
% BEGIN S2.SS0.SSS0.Px1.p7
 
Higher-order $\beta$-matching is the following typed unification problem, for which only one side is subject to instantiation.\par
 
% END S2.SS0.SSS0.Px1.p7
 
% BEGIN Thmtheorem7
 
\par\medskip\noindent\textbf{Problem 7 (Higher-order $\beta$-Matching (${F}\,\mathsf{X}=N$)).}\par
 
% BEGIN Thmtheorem7.p1
 
Given terms $F,N$ and simple types $\sigma,\tau$ such that $\emptyset\vdash F:\sigma\to\tau$ and $\emptyset\vdash N:\tau$, is there a term $M$ such that $\emptyset\vdash M:\sigma$ and $F\,M=_{\beta}N$?\par
 
% END Thmtheorem7.p1
 
% END Thmtheorem7
 
% BEGIN S2.SS0.SSS0.Px1.p8
 
Undecidability of higher-order $\beta$-matching is shown by Loader [13] using a reduction from a variant of $\lambda$-definability.\par
 
% END S2.SS0.SSS0.Px1.p8
 
% BEGIN Thmtheorem8
 
\par\medskip\noindent\textbf{Theorem 8 ([13, Theorem 5.5]).}\par
 
% BEGIN Thmtheorem8.p1
 
Higher-order $\beta$-matching is undecidable.\par
 
% END Thmtheorem8.p1
 
% END Thmtheorem8
 
% BEGIN S2.SS0.SSS0.Px1.p9
 
For the remainder of the present work we use the term matching in order to refer to higher-order $\beta$-matching. Since simply typed terms are strongly normalizing and $\beta$-reduction is confluent [2], the it suffices in Problem 7 to consider terms $F,M,N$ in normal form.\par
 
% END S2.SS0.SSS0.Px1.p9
 
% BEGIN S2.SS0.SSS0.Px1.p10
 
Let us get familiar with matching by means of several illustrating examples. The following Example 9 illustrates a positive matching instance.\par
 
% END S2.SS0.SSS0.Px1.p10
 
% BEGIN Thmtheorem9
 
\par\medskip\noindent\textbf{Example 9.}\par
 
% BEGIN Thmtheorem9.p1
 
Consider the terms $F:=\lambda x.\lambda y.x\,y\,I$ and $N:=I$, for which we have $\emptyset\vdash F:(\iota\to(\iota\to\iota)\to\iota)\to(\iota\to\iota)$ and $\emptyset\vdash N:\iota\to\iota$.\par
 
% END Thmtheorem9.p1
 
% BEGIN Thmtheorem9.p2
 
The matching instance ${F}\,\mathsf{X}=N$ is solvable, including the solution $M:=\lambda u.\lambda f.f\,u$.To be precise, we have the following two properties:\par
  
\par\noindent  $\blacksquare$ 
% BEGIN S2.I1.i1.p1
 
$\emptyset\vdash M:\iota\to(\iota\to\iota)\to\iota$ (Example 6)\par
 
% END S2.I1.i1.p1
 
 
\par\noindent  $\blacksquare$ 
% BEGIN S2.I1.i2.p1
 
$F\,M=_{\beta}\lambda y.(\lambda u.\lambda f.f\,u)\,y\,I=_{\beta}\lambda y.I\,y%
=_{\beta}N$\par
 
% END S2.I1.i2.p1
 
  
% END Thmtheorem9.p2
 
% END Thmtheorem9
 
% BEGIN S2.SS0.SSS0.Px1.p11
 
We want to encode certain functional behavior as a matching instance. The following Example 10 shows a naive approach to such an encoding and its limitations.\par
 
% END S2.SS0.SSS0.Px1.p11
 
% BEGIN Thmtheorem10
 
\par\medskip\noindent\textbf{Example 10.}\par
 
% BEGIN Thmtheorem10.p1
 
Let us associate elements of the set $\{1,2,3\}$ with projections $\pi_{1}:=\lambda xyz.x$, $\pi_{2}:=\lambda xyz.y$, and $\pi_{3}:=\lambda xyz.z$ respectively. For the term $G:=\lambda h.\lambda xyz.h\,y\,z\,x$ we have $G\,\pi_{1}=_{\beta}\pi_{2}$, $G\,\pi_{2}=_{\beta}\pi_{3}$, and $G\,\pi_{3}=_{\beta}\pi_{1}$. Therefore, $G$ realizes a finite function $f_{G}:\{1,2,3\}\to\{1,2,3\}$ such that $f_{G}(1)=2$, $f_{G}(2)=3$, and $f_{G}(3)=1$.\par
 
% END Thmtheorem10.p1
 
% BEGIN Thmtheorem10.p2
 
Let $\kappa:=\iota\to\iota\to\iota\to\iota$ be a simple type for which we have $\emptyset\vdash\pi_{i}:\kappa$ for $i\in\{1,2,3\}$. Consider the matching instance $F\,\mathsf{X}=\pi_{3}$, where $F:=\lambda t.t\,G\,\pi_{1}$, and for which we have $\emptyset\vdash F:((\kappa\to\kappa)\to\kappa\to\kappa)\to\kappa$. The intended “meaning” of this matching instance is: starting with the element $1$, repeatedly apply the function $f_{G}$ in order to construct the element $3$.\par
 
% END Thmtheorem10.p2
 
% BEGIN Thmtheorem10.p3
 
One solution for this instance is the term $\lambda f.\lambda s.f\,(f\,s)$ for which we have:\par
    \[F\,(\lambda f.\lambda s.f\,(f\,s))=_{\beta}f\,(f\,s)[f:=G,s:=\pi_{1}]=_{\beta}%
G\,\pi_{2}=_{\beta}\pi_{3}\]    
This solution follows the intended meaning of the underlying representation, constructing the element $f_{G}(f_{G}(1))=3$. Another solution is $\lambda f.\lambda s.f\,(f\,(f\,(f\,(f\,s))))$, which utilizes $f_{G}(f_{G}(f_{G}(f_{G}(f_{G}(1)))))=3$.\par
 
% END Thmtheorem10.p3
 
% BEGIN Thmtheorem10.p4
 
Unfortunately, there are solutions to the above matching instance which behave differently. One such solution is $\lambda f.\lambda s.\pi_{3}$, for which we also have $F\,(\lambda f.\lambda s.\pi_{3})=_{\beta}\pi_{3}$. In this case, the element $3$ is constructed “ad hoc”, with no reference to the provided arguments. Another solution is the term $\lambda f.\lambda s.\lambda xyz.f\,s\,z\,z\,z$. This solution exploits the exact representation of elements via projections, disregarding any intended meaning of the underlying representation.\par
 
% END Thmtheorem10.p4
 
% END Thmtheorem10
 
% BEGIN S2.SS0.SSS0.Px1.p12
 
In the above Example 10 we would like to exclude certain ad hoc solutions, in order to faithfully encode intended functional behavior. Towards this aim, a known technique how to restrict the shape of solutions is illustrated in the following Example 11.\par
 
% END S2.SS0.SSS0.Px1.p12
 
% BEGIN Thmtheorem11
 
\par\medskip\noindent\textbf{Example 11 ([5, Proposition 3.4]).}\par
 
% BEGIN Thmtheorem11.p1
 
Consider a term $M$ in normal form such that $M\,I\,u=_{\beta}u$ where $u$ is a term variable, and $\emptyset\vdash M:(\kappa\to\kappa)\to\kappa\to\kappa$. By case analysis on $M$ we have that $M=\lambda f.\lambda s.N$ for some term $N$ in normal form. Furthermore, $\{f:\kappa\to\kappa,s:\kappa\}\vdash N:\kappa$ and $N[f:=\lambda x.x,s:=u]=_{\beta}u$. Therefore, the term $N$ is not an abstraction. By induction on the size of $N$ and case analysis of the normal form we have that $N=s$ or $N=f\,(\ldots(f\,s)\ldots)$. Since the term $M$ contains exactly two $\lambda$-abstractions, it cannot be an ad hoc solution from Example 10.\par
 
% END Thmtheorem11.p1
 
% END Thmtheorem11
 
% BEGIN S2.SS0.SSS0.Px1.p13
 
Combining Example 10 with Example 11, we formulate in the following Example 12 a matching instance which faithfully encodes the desired functional behavior.\par
 
% END S2.SS0.SSS0.Px1.p13
 
% BEGIN Thmtheorem12
 
\par\medskip\noindent\textbf{Example 12.}\par
 
% BEGIN Thmtheorem12.p1
 
Let $F:=\lambda t.\lambda r.r\,(t\,G\,\pi_{1})\,(\lambda u.t\,I\,u)$ and $N:=\lambda r.r\,\pi_{3}\,(\lambda u.u)$ where the term $G=\lambda h.\lambda xyz.h\,y\,z\,x$ is from Example 10. We have\par
  
\par\noindent  $\blacksquare$ 
% BEGIN S2.I2.i1.p1
 
$\emptyset\vdash F:((\kappa\to\kappa)\to\kappa\to\kappa)\to(\kappa\to(\kappa\to%
\kappa)\to\iota)\to\iota$\par
 
% END S2.I2.i1.p1
 
 
\par\noindent  $\blacksquare$ 
% BEGIN S2.I2.i2.p1
 
$\emptyset\vdash N:(\kappa\to(\kappa\to\kappa)\to\iota)\to\iota$\par
 
% END S2.I2.i2.p1
 
  
The matching instance $F\,\mathsf{X}=N$ combines the matching instance from Example 10 with the additional restriction from Example 11. Therefore, solutions such as $\lambda f.\lambda s.f\,(f\,s)$ and $\lambda f.\lambda s.f\,(f\,(f\,(f\,(f\,s))))$ from Example 10 which follow the intended “meaning” of the underlying representation still solve $F\,\mathsf{X}=N$. However, ad hoc solutions such as $\lambda f.\lambda s.\pi_{3}$ or $\lambda f.\lambda s.\lambda xyz.f\,s\,z\,z\,z$ from Example 10 do not solve $F\,\mathsf{X}=N$ because such solutions contain too many $\lambda$-abstractions.\par
 
% END Thmtheorem12.p1
 
% END Thmtheorem12
 
% BEGIN Thmtheorem13
 
\par\medskip\noindent\textbf{$\blacktriangleright$ Remark 13.}\par
 
% BEGIN Thmtheorem13.p1
 
Without the restriction $\emptyset\vdash M:(\kappa\to\kappa)\to\kappa\to\kappa$ in Example 11 the term $M:=\lambda f.\lambda s.f\,f\,s$ satisfies $M\,I\,u=_{\beta}u$. However, Example 12 does not admit $M$ as a solution. The present work relies on well-typedness, but might be adapted to an untyped scenario.\par
 
% END Thmtheorem13.p1
 
% END Thmtheorem13
 
% BEGIN S2.SS0.SSS0.Px1.p14
 
The following Remark 14 illustrates how the addition of $\eta$-reduction would make the technique from Example 12 (as well as Loader’s approach) unsuitable.\par
 
% END S2.SS0.SSS0.Px1.p14
 
% BEGIN Thmtheorem14
 
\par\medskip\noindent\textbf{$\blacktriangleright$ Remark 14.}\par
 
% BEGIN Thmtheorem14.p1
 
Example 11 demonstrates how to restrict the number of abstractions in solutions. However, in the presence of $\eta$-reduction (contextual closure of $\lambda x.f\,x\to_{\eta}f$) this does not work, as shown below. Consider the terms $G:=\lambda h.\lambda xyz.h\,y\,z\,x$, $\pi_{1}=\lambda xyz.x$, $\pi_{2}=\lambda xyz.y$, and $\pi_{3}=\lambda xyz.z$ from Example 11. The term $M:=\lambda g.\lambda h.\lambda xyz.h\,(g\,\pi_{1}x\,z\,y)\,y\,z$ solves the matching instance in Example 12 because:\par
    \[\begin{array}[t]{lclcl}M\,G\,\pi_{1}&=_{\beta}&\big{(}\lambda g.\lambda h.%
\lambda xyz.h\,(g\,\pi_{1}x\,z\,y)\,y\,z\big{)}\,G\,\pi_{1}&=_{\beta}&\lambda
xyz%
.\pi_{1}\,(G\,\pi_{1}x\,z\,y)\,y\,z\\
&=_{\beta}&\lambda xyz.G\,\pi_{1}x\,z\,y&=_{\beta}&\lambda xyz.\pi_{1}\,z\,y\,%
x=_{\beta}\lambda xyz.z=\pi_{3}\end{array}\]       \[\begin{array}[t]{lclcl}M\,I\,u&=_{\beta}&\big{(}\lambda g.\lambda h.\lambda xyz%
.h\,(g\,\pi_{1}x\,z\,y)\,y\,z\big{)}\,I\,u&=_{\beta}&\lambda xyz.u\,(I\,\pi_{1%
}x\,z\,y)\,y\,z\\
&=_{\beta}&\lambda xyz.u\,(\pi_{1}x\,z\,y)\,y\,z&=_{\beta}&\lambda xyz.u\,x\,y%
\,z\to_{\eta}^{*}u\end{array}\]    
In particular, $\eta$-reduction allows for additional $\lambda$-abstractions in the solution, making the technique from Example 11 unsuitable.\par
 
% END Thmtheorem14.p1
 
% END Thmtheorem14
 
% BEGIN S2.SS0.SSS0.Px1.p15
 
The observation from Example 12 is generalized by Loader to encode arbitrary families of finite functions. This results in undecidability of higher-order $\beta$-matching by reduction from a variant of $\lambda$-definability. Loader’s generalization is quite sophisticated, as it requires construction principles to restrict shapes of realizers of arbitrary higher-order finite functions. In the present work, we focus on a fragment which can be identified by inspection of intersection types occurring in the undecidability proof of intersection type inhabitation [20, 6] and their relationship to finite model theory [15]. This leads to a simpler undecidability proof and reveals a connection between matching, intersection type inhabitation, and $\lambda$-definability. The presented approach has similarities with Joly’s [11] refinement of Loader’s proof. Joly shifts the technical burden to a particular $\lambda$-definability problem, which then is simpler to handle. Instead, we avoid $\lambda$-definability altogether and use a rewriting problem in a restricted class of simple semi-Thue systems as a starting point.\par
 
% END S2.SS0.SSS0.Px1.p15
 
% END S2.SS0.SSS0.Px1
 
% BEGIN S2.SS0.SSS0.Px2
 
\par\medskip\noindent\textbf{Simple Semi-Thue Systems}\par
 
% BEGIN S2.SS0.SSS0.Px2.p1
 
A simple semi-Thue system (Definition 15) is a rewriting system of restricted shape, introduced by Urzyczyn [20] in order to show undecidability of intersection type inhabitation.\par
 
% END S2.SS0.SSS0.Px2.p1
 
% BEGIN Thmtheorem15
 
\par\medskip\noindent\textbf{Definition 15 (Simple Semi-Thue System).}\par
 
% BEGIN Thmtheorem15.p1
 
A semi-Thue system $\mathfrak{R}$ over an alphabet $\mathcal{A}$ is simple, if each rule has the shape $ab\Rightarrow cd$ for $a,b,c,d\in\mathcal{A}$.\par
 
% END Thmtheorem15.p1
 
% END Thmtheorem15
 
% BEGIN S2.SS0.SSS0.Px2.p2
 
The reflexive, transitive closure of the rewriting relation for a given simple semi-Thue system $\mathfrak{R}$ is denoted $\Rightarrow_{\mathfrak{R}}^{*}$. For arbitrary simple semi-Thue systems it is undecidable whether some non-empty sequence of $\mathbf{0}$s can be transformed into a non-empty sequence of $\mathbf{1}$s.\par
 
% END S2.SS0.SSS0.Px2.p2
 
% BEGIN Thmtheorem16
 
\par\medskip\noindent\textbf{Problem 16 ($\mathbf{0}^{+}\Rightarrow^{*}\mathbf{1}^{+}$).}\par
 
% BEGIN Thmtheorem16.p1
 
Given a simple semi-Thue system $\mathfrak{R}$, does $\mathbf{0}^{n}\Rightarrow^{*}_{\mathfrak{R}}\mathbf{1}^{n}$ hold for some $n>0$?\par
 
% END Thmtheorem16.p1
 
% END Thmtheorem16
 
% BEGIN Thmtheorem17
 
\par\medskip\noindent\textbf{Theorem 17 ([6, Lemma 3.3]).}\par
 
% BEGIN Thmtheorem17.p1
 
Problem $\mathbf{0}^{+}\Rightarrow^{*}\mathbf{1}^{+}$ is undecidable.\par
 
% END Thmtheorem17.p1
 
% END Thmtheorem17
 
% BEGIN S2.SS0.SSS0.Px2.p3
 
The following Example 18 illustrates a positive instance of Problem $\mathbf{0}^{+}\Rightarrow^{*}\mathbf{1}^{+}$.\par
 
% END S2.SS0.SSS0.Px2.p3
 
% BEGIN Thmtheorem18
 
\par\medskip\noindent\textbf{Example 18.}\par
 
% BEGIN Thmtheorem18.p1
 
Let $\mathfrak{R}:=\{\mathbf{0}\mathbf{0}\Rightarrow\mathbf{2}\mathbf{2},\mathbf{0}%
\mathbf{2}\Rightarrow\mathbf{1}\mathbf{1},\mathbf{2}\mathbf{0}\Rightarrow%
\mathbf{1}\mathbf{1}\}$ be a simple semi-Thue system over the alphabet $\{\mathbf{0},\mathbf{1},\mathbf{2}\}$. We have $\mathbf{0}\mathbf{0}\mathbf{0}\mathbf{0}\Rightarrow_{\mathfrak{R}}\mathbf{0}%
\mathbf{2}\mathbf{2}\mathbf{0}\Rightarrow_{\mathfrak{R}}\mathbf{1}\mathbf{1}%
\mathbf{2}\mathbf{0}\Rightarrow_{\mathfrak{R}}\mathbf{1}\mathbf{1}\mathbf{1}%
\mathbf{1}$. As a side note, we have $\mathbf{0}^{n}\;{\not\Rightarrow}_{\mathfrak{R}}^{*}\;\mathbf{1}^{n}$ for $n\in\{1,2,3\}$.\par
 
% END Thmtheorem18.p1
 
% END Thmtheorem18
 
% BEGIN Thmtheorem19
 
\par\medskip\noindent\textbf{$\blacktriangleright$ Remark 19.}\par
 
% BEGIN Thmtheorem19.p1
 
Problem $\mathbf{0}^{+}\Rightarrow^{*}\mathbf{1}^{+}$ is used as a starting point in a refinement [6, Lemma 4.4] of Urzyczyn’s undecidability result for intersection type inhabitation [20]. Undecidability of Problem $\mathbf{0}^{+}\Rightarrow^{*}\mathbf{1}^{+}$ is mechanized as part of Coq Library of Undecidability Proofs [8], making it a good starting point for further mechanized undecidability results.\par
 
% END Thmtheorem19.p1
 
% END Thmtheorem19
 
% END S2.SS0.SSS0.Px2
 
% END S2


% BEGIN S3
 
\par\medskip\noindent\textbf{3 Undecidability of Higher-order $\beta$-Matching}\par
 
% BEGIN S3.p1
 
In this section we develop our main result (Theorem 39): a reduction from the rewriting problem $\mathbf{0}^{+}\Rightarrow^{*}\mathbf{1}^{+}$ to higher-order $\beta$-matching.\par
 
% END S3.p1
 
% BEGIN S3.p2
 
For the remainder of the section we fix the simple semi-Thue system $\mathfrak{R}:=\{R_{1},\ldots,R_{L}\}$ with $L>0$ rules over the finite alphabet $\{\mathbf{0},\mathbf{1},\mathbf{2},\ldots,\mathbf{K}\}$. Our approach is to construct simply typed terms which capture the two main aspects of the rewriting problem $\mathbf{0}^{+}\Rightarrow^{*}\mathbf{1}^{+}$: the search for a sufficiently long starting sequence of $\mathbf{0}$s, and the individual rewriting steps to the desired sequence of $\mathbf{1}$s.\par
 
% END S3.p2
 
% BEGIN S3.p3
 
The remainder of the present section is structured as follows. First, we fix basic notation, encoding, and properties of the rewriting in the system $\mathfrak{R}$. Second, we restrict the shape of potential solutions for the constructed matching instance, similarly to Example 11. Third, for solutions of restricted shape we capture the functional properties of the rewriting problem $\mathbf{0}^{+}\Rightarrow^{*}\mathbf{1}^{+}$.\par
 
% END S3.p3
 
% BEGIN S3.SSx1
 
\par\medskip\noindent\textbf{Notation}\par
 
% BEGIN S3.SSx1.p1
 
We introduce four additional symbols in extended alphabet $\mathcal{A}:=\{\mathbf{0},\mathbf{1},\ldots,\mathbf{K}\}\cup\{\$,\bullet,\top,\bot\}$. We represent an alphabet symbol $i\in\mathcal{A}$ as the projection $\pi_{i}:=\lambda s_{\mathbf{0}}s_{\mathbf{1}}\ldots s_{\mathbf{K}}s_{\$}s_{%
\bullet}s_{\top}s_{\bot}.s_{i}$ typed by the simple type $\kappa:=\underbracket{\iota\to\ldots\to\iota}_{|\mathcal{A}|\leavevmode%
\nobreak\ \text{times}}\to\iota$. The intended meaning of the additional symbols in the extended alphabet is as follows. The symbol $\$$ marks the beginning of a word for word expansion, $\bullet$ serves as symbol different from $\mathbf{0}$ and $\mathbf{1}$ for case analysis, $\top$ serves as a constructible symbol, and $\bot$ serves as a non-constructible symbol.\par
 
% END S3.SSx1.p1
 
% BEGIN S3.SSx1.p2
 
For readability, we use the following case notation to match individual symbols.\par
 
% END S3.SSx1.p2
 
% BEGIN Thmtheorem20
 
\par\medskip\noindent\textbf{Definition 20 (case).}\par
 
% BEGIN Thmtheorem20.p1
 
For $k\in\mathbb{N}$, distinct $i_{1},\ldots,i_{k}\in\mathcal{A}$, and terms $M_{1},\ldots,M_{k}$:\par
    $\displaystyle\textnormal{{case }}{x}\textnormal{{ of }}\langle M\mid i_{1}%
\mapsto M_{1}\mid\ldots\mid i_{k}\mapsto M_{k}\rangle:=$ $\displaystyle\leavevmode\nobreak\ x\,N_{\mathbf{0}}\,N_{\mathbf{1}}\ldots N_{%
\mathbf{K}}\,N_{\$}\,N_{\bullet}\,N_{\top}\,N_{\bot}$      $\displaystyle\text{ where }N_{i}=\begin{cases}M_{j}&\text{if }i=i_{j}\\
M&\text{otherwise}\end{cases}$    
% END Thmtheorem20.p1
 
% END Thmtheorem20
 
% BEGIN S3.SSx1.p3
 
A particular term $\delta_{i}$ for $i\in\mathcal{A}$, which we use commonly is:\par
    \[\delta_{i}:=\lambda x.\lambda s_{\mathbf{0}}s_{\mathbf{1}}\ldots s_{\mathbf{K}%
}s_{\$}s_{\bullet}s_{\top}s_{\bot}.\textnormal{{case }}{x}\textnormal{{ of }}%
\langle s_{\bot}\mid\top\mapsto s_{i}\rangle\]    
We have $\emptyset\vdash\delta_{i}:\kappa\to\kappa$, and the following Lemma 21 specifies the behavior of $\delta_{i}$.\par
 
% END S3.SSx1.p3
 
% BEGIN Thmtheorem21
 
\par\medskip\noindent\textbf{Lemma 21.}\par
 
% BEGIN Thmtheorem21.p1
 
For $i,j\in\mathcal{A}$ such that $i\neq j$ we have $\delta_{i}\,\pi_{\top}=_{\beta}\pi_{i}$ and $\delta_{i}\,\pi_{j}=_{\beta}\pi_{\bot}$.\par
 
% END Thmtheorem21.p1
 
% END Thmtheorem21
 
% END S3.SSx1
 
% BEGIN S3.SSx2
 
\par\medskip\noindent\textbf{Syntactic Constraints}\par
 
% BEGIN S3.SSx2.p1
 
We identify the shape of “well-formed” terms, suitable to represent rewriting. In the following Definition 22 terms in the set $\mathcal{Q}_{m}$ capture consecutive rule application for a word of length $m+1$, ending in the word $\mathbf{1}^{m+1}$ (represented by $z_{\mathbf{1}}\in\mathcal{Q}_{m}$). The subterm $r_{i}\,p_{j}$ (and $r_{i}\,(\lambda w.p_{j}\,w)$) for $i\in\{1,\ldots,L\}$ and $j\in\{1,\ldots,m\}$ indicates an application of the rule $R_{i}$ at position $j$. Additionally, terms in the set $\mathcal{R}_{m}$ capture consecutive increase of word length starting with $m+1$, and initialization with $\mathbf{0}$s before rewriting (represented by $z_{\mathbf{0}}\,N\,M\in\mathcal{R}_{m}$ for $M\in\mathcal{Q}_{m}$). Specifically, the subterm $z_{\star}\,N\,(\lambda p_{m+1}.M)$ introduces an additional bound variable $p_{m+1}$ in order to argue about rule application at position $m+1$ in the longer word. Consequently, terms in $\mathcal{R}_{1}$ represent witnesses for an arbitrary expansion of a word starting with length $2$, followed by initialization with $\mathbf{0}$s, and consecutive rule application (potentially ending in $\mathbf{1}$s).\par
 
% END S3.SSx2.p1
 
% BEGIN Thmtheorem22
 
\par\medskip\noindent\textbf{Definition 22 (Sets $\mathcal{Q}_{m}$, $\mathcal{R}_{m}$ of Terms).}\par
 
% BEGIN Thmtheorem22.p1
 
For $m>0$, let $\mathcal{Q}_{m}$ and $\mathcal{R}_{m}$ be the smallest sets of terms satisfying the following rules:\par
  
\par\noindent  $\blacksquare$ 
% BEGIN S3.I1.i1.p1
 
$z_{\mathbf{1}}\in\mathcal{Q}_{m}$\par
 
% END S3.I1.i1.p1
 
 
\par\noindent  $\blacksquare$ 
% BEGIN S3.I1.i2.p1
 
if $M\in\mathcal{Q}_{m}$ then $(r_{i}\,p_{j}\,M)\in\mathcal{Q}_{m}$ for $i\in\{1,\ldots,L\}$ and $j\in\{1,\ldots,m\}$\par
 
% END S3.I1.i2.p1
 
 
\par\noindent  $\blacksquare$ 
% BEGIN S3.I1.i3.p1
 
if $M\in\mathcal{Q}_{m}$ then $(r_{i}\,(\lambda w.p_{j}\,w)\,M)\in\mathcal{Q}_{m}$ for $i\in\{1,\ldots,L\}$ and $j\in\{1,\ldots,m\}$\par
 
% END S3.I1.i3.p1
 
 
\par\noindent  $\blacksquare$ 
% BEGIN S3.I1.i4.p1
 
if $M\in\mathcal{Q}_{m}$ then $(z_{\mathbf{0}}\,N\,M)\in\mathcal{R}_{m}$\par
 
% END S3.I1.i4.p1
 
 
\par\noindent  $\blacksquare$ 
% BEGIN S3.I1.i5.p1
 
if $M\in\mathcal{R}_{m+1}$ then $(z_{\star}\,N\,(\lambda p_{m+1}.M))\in\mathcal{R}_{m}$\par
 
% END S3.I1.i5.p1
 
  
% END Thmtheorem22.p1
 
% END Thmtheorem22
 
% BEGIN Thmtheorem23
 
\par\medskip\noindent\textbf{$\blacktriangleright$ Remark 23.}\par
 
% BEGIN Thmtheorem23.p1
 
Terms in $\mathcal{R}_{1}$ are inspired by inhabitants in a refinement [6, Lemma 4.4] of Urzyczyn’s undecidability result for intersection type inhabitation [20].\par
 
% END Thmtheorem23.p1
 
% END Thmtheorem23
 
% BEGIN S3.SSx2.p2
 
Free variables occurring in terms in $\mathcal{Q}_{m}$ and $\mathcal{R}_{m}$ are assigned simple types according to the following type environment $\Gamma_{m}$, such that terms in $\mathcal{Q}_{m}$ and $\mathcal{R}_{m}$ can be assigned the simple type $\kappa$.\par
 
% END S3.SSx2.p2
 
% BEGIN Thmtheorem24
 
\par\medskip\noindent\textbf{Definition 24 (Type Environment $\Gamma_{m}$).}\par
 
% BEGIN Thmtheorem24.p1
 
For $m>0$ let\par
    $\displaystyle\Gamma_{m}$ $\displaystyle:=\{\begin{aligned} &z_{\mathbf{1}}:\kappa,z_{\mathbf{0}}:(\kappa%
\to\kappa)\to\kappa\to\kappa,z_{\star}:(\kappa\to\kappa)\to((\kappa\to\kappa)%
\to\kappa)\to\kappa,\\
&p_{1}:\kappa\to\kappa,\ldots,p_{m}:\kappa\to\kappa,\\
&r_{1}:(\kappa\to\kappa)\to\kappa\to\kappa,\ldots,r_{L}:(\kappa\to\kappa)\to%
\kappa\to\kappa\}\end{aligned}$    
% END Thmtheorem24.p1
 
% END Thmtheorem24
 
% BEGIN S3.SSx2.p3
 
Similarly to Example 11, we formulate typed terms (Definition 25) and $\beta$-equivalence constraints characterizing members of $\mathcal{Q}_{m}$ (Lemma 27) and $\mathcal{R}_{m}$ (Lemma 28).\par
 
% END S3.SSx2.p3
 
% BEGIN Thmtheorem25
 
\par\medskip\noindent\textbf{Definition 25 (Typed Terms $H_{\star}$, $H_{\mathbf{0}}$, $H_{R}$).}\par
 
% BEGIN Thmtheorem25.p1
    $\displaystyle H_{\star}:=$ $\displaystyle\lambda h.\lambda g.\lambda s_{\mathbf{0}}s_{\mathbf{1}}\ldots s_%
{\mathbf{K}}s_{\$}s_{\bullet}s_{\top}s_{\bot}.\textnormal{{case }}{g\,\delta_{%
\bullet}}\textnormal{{ of }}\langle s_{\bot}\mid\$\mapsto s_{\$}\rangle$     $\displaystyle H_{\mathbf{0}}:=$ $\displaystyle\lambda h.\lambda x.\lambda s_{\mathbf{0}}s_{\mathbf{1}}\ldots s_%
{\mathbf{K}}s_{\$}s_{\bullet}s_{\top}s_{\bot}.\textnormal{{case }}{x}%
\textnormal{{ of }}\langle s_{\bot}\mid\mathbf{1}\mapsto s_{\$}\rangle$     $\displaystyle H_{R}:=$ $\displaystyle\lambda h.\lambda x.\lambda s_{\mathbf{0}}s_{\mathbf{1}}\ldots s_%
{\mathbf{K}}s_{\$}s_{\bullet}s_{\top}s_{\bot}.\textnormal{{case }}{h\,\pi_{%
\top}}\textnormal{{ of }}\langle s_{\bot}\mid\bullet\mapsto\textnormal{{case }%
}{x}\textnormal{{ of }}\langle s_{\bot}\mid\mathbf{1}\mapsto s_{\mathbf{1}}\rangle\rangle$       \[\emptyset\vdash H_{\star}:\Gamma_{m}(z_{\star})\qquad\quad\emptyset\vdash H_{%
\mathbf{0}}:\Gamma_{m}(z_{\mathbf{0}})\qquad\quad\emptyset\vdash H_{R}:\Gamma_%
{m}(r_{i})\text{ for }i\in\{1,\ldots,L\}\]    
% END Thmtheorem25.p1
 
% END Thmtheorem25
 
% BEGIN S3.SSx2.p4
 
We introduce substitutions $S_{F}$ and $S_{H}$ acting on the term variables $z_{\star},z_{\mathbf{1}},z_{\mathbf{0}},r_{1},\ldots,r_{L}$, which occur in terms in $\mathcal{Q}_{m}$ and $\mathcal{R}_{m}$.\par
 
% END S3.SSx2.p4
 
% BEGIN Thmtheorem26
 
\par\medskip\noindent\textbf{Definition 26 (Substitutions $S_{F}$, $S_{H}$).}\par
 
% BEGIN Thmtheorem26.p1
    \[\begin{array}[]{llll}S_{F}(z_{\star}):=\lambda h.\lambda g.g\,I\nobreak%
\leavevmode&S_{F}(z_{\mathbf{1}}):=u\nobreak\leavevmode&S_{F}(z_{\mathbf{0}}):%
=\lambda h.I\nobreak\leavevmode&S_{F}(r_{j}):=I\text{ for }j\in{1,\ldots L}\\
S_{H}(z_{\star}):=H_{\star}&S_{H}(z_{\mathbf{1}}):=\pi_{\mathbf{1}}&S_{H}(z_{%
\mathbf{0}}):=H_{\mathbf{0}}&S_{H}(r_{j}):=H_{R}\text{ for }j\in{1,\ldots L}%
\end{array}\]    
% END Thmtheorem26.p1
 
% END Thmtheorem26
 
% BEGIN Thmtheorem27
 
\par\medskip\noindent\textbf{Lemma 27.}\par
 
% BEGIN Thmtheorem27.p1
 
For $m>0$, if a term $M$ is in normal form such that $\Gamma_{m}\vdash M:\kappa$, $S_{F}(M)[p_{1}{\,:=\,}I,\ldots,p_{m}{\,:=\,}I]=_{\beta}u$, and $S_{H}(M)[p_{1}:=\delta_{\bullet},\ldots,p_{m}:=\delta_{\bullet}]=_{\beta}\pi_{%
\mathbf{1}}$, then $M\in\mathcal{Q}_{m}$.\par
 
% END Thmtheorem27.p1
 
% END Thmtheorem27
 
% BEGIN S3.SSx2.SSS0.P0.SPx1
 
\par\medskip\noindent\textbf{Proof.}\par
 
% BEGIN S3.SSx2.SSS0.P0.SPx1.p1
 
Induction on the size of $M$ and case analysis of the normal form. $\hfill\blacktriangleleft$\par
 
% END S3.SSx2.SSS0.P0.SPx1.p1
 
% BEGIN Thmtheorem28
 
\par\medskip\noindent\textbf{Lemma 28.}\par
 
% BEGIN Thmtheorem28.p1
 
For $m>0$, if a term $M$ is in normal form such that $\Gamma_{m}\vdash M:\kappa$, $S_{F}(M)[p_{1}{\,:=\,}I,\ldots,p_{m}{\,:=\,}I]=_{\beta}u$, and $S_{H}(M)[p_{1}:=\delta_{\bullet},\ldots,p_{m}:=\delta_{\bullet}]=_{\beta}\pi_{\$}$, then $M\in\mathcal{R}_{m}$.\par
 
% END Thmtheorem28.p1
 
% END Thmtheorem28
 
% END S3.SSx2.SSS0.P0.SPx1
 
% BEGIN S3.SSx2.SSS0.P0.SPx2
 
\par\medskip\noindent\textbf{Proof.}\par
 
% BEGIN S3.SSx2.SSS0.P0.SPx2.p1
 
Induction on the size of $M$, case analysis of the normal form, and Lemma 27. $\hfill\blacktriangleleft$\par
 
% END S3.SSx2.SSS0.P0.SPx2.p1
 
% BEGIN S3.SSx2.SSS0.P0.SPx2.p2
 
As a consequence of the above Lemma 28, the following Theorem 29 presents $\beta$-equivalence constraints which suffice to restrict the shape of terms under consideration.\par
 
% END S3.SSx2.SSS0.P0.SPx2.p2
 
% BEGIN Thmtheorem29
 
\par\medskip\noindent\textbf{Theorem 29.}\par
 
% BEGIN Thmtheorem29.p1
 
If a term $M$ is in normal form such that $\emptyset\vdash M:\Gamma_{1}(r_{1})\to\cdots\to\Gamma_{1}(r_{L})\to\Gamma_{1}(%
z_{\mathbf{0}})\to\Gamma_{1}(z_{\mathbf{1}})\to\Gamma_{1}(z_{\star})\to\Gamma_%
{1}(p_{1})\to\kappa$, $M\,I\ldots I\,(\lambda h.I)\,u\,(\lambda h.\lambda g.g\,I)\,I=_{\beta}u$, and $M\,H_{R}\ldots H_{R}\,H_{\mathbf{0}}\,\pi_{\mathbf{1}}\,H_{\star}\,\delta_{%
\bullet}=_{\beta}\pi_{\$}$, then $M=\lambda r_{1}\ldots r_{L}.\lambda z_{\mathbf{0}}z_{\mathbf{1}}z_{\star}p_{1}.N$ for some term $N$ such that $N\in\mathcal{R}_{1}$.\par
 
% END Thmtheorem29.p1
 
% END Thmtheorem29
 
% END S3.SSx2.SSS0.P0.SPx2
 
% BEGIN S3.SSx2.SSS0.P0.SPx3
 
\par\medskip\noindent\textbf{Proof.}\par
 
% BEGIN S3.SSx2.SSS0.P0.SPx3.p1
 
Induction on the size of $M$, case analysis of the normal form, and Lemma 28. $\hfill\blacktriangleleft$\par
 
% END S3.SSx2.SSS0.P0.SPx3.p1
 
% BEGIN Thmtheorem30
 
\par\medskip\noindent\textbf{Example 30.}\par
 
% BEGIN Thmtheorem30.p1
 
Assume $\mathfrak{R}=\{\mathbf{0}\mathbf{0}\Rightarrow\mathbf{2}\mathbf{2},\mathbf{0}%
\mathbf{2}\Rightarrow\mathbf{1}\mathbf{1},\mathbf{2}\mathbf{0}\Rightarrow%
\mathbf{1}\mathbf{1}\}$ over the alphabet $\{\mathbf{0},\mathbf{1},\mathbf{2}\}$ from Example 18. Let $N:=r_{1}\,p_{2}\,(r_{2}\,p_{1}\,(r_{3}\,p_{3}\,z_{\mathbf{1}}))$ and $M:=z_{\star}\,p_{1}\,(\lambda p_{2}.z_{\star}\,p_{2}\,(\lambda p_{3}.z_{%
\mathbf{0}}\,p_{3}\,N))$. We have $N\in\mathcal{Q}_{3}$ and $M\in\mathcal{R}_{1}$. In congruence with Theorem 29 we have:\par
    \[\begin{array}[]{ll}(\lambda r_{1}r_{2}r_{3}.\lambda z_{\mathbf{0}}z_{\mathbf{1%
}}z_{\star}p_{1}.M)\,I\,I\,I\,(\lambda h.I)\,u\,(\lambda h.\lambda g.g\,I)\,I&%
=_{\beta}u\\
(\lambda r_{1}r_{2}r_{3}.\lambda z_{\mathbf{0}}z_{\mathbf{1}}z_{\star}p_{1}.M)%
\,H_{R}\,H_{R}\,H_{R}\,H_{\mathbf{0}}\,\pi_{\mathbf{1}}\,H_{\star}\,\delta_{%
\bullet}&=_{\beta}\pi_{\$}\end{array}\]    
While Theorem 29 only establishes “well-formedness”, the term $M$ has an intended meaning: An initial word of length $2$ is expanded twice (using $z_{\star}$) to a word of length $4$ and initialized to $\mathbf{0}$s (using $z_{\mathbf{0}}$). The introduced variables $p_{2}$ and $p_{3}$ are used to address positions in the longer word. The intended meaning of $N$ is that the first rule (using $r_{1}$) is applied at position $2$ (using $p_{2}$), followed by the second rule at position $1$, and third rule at position $3$ accordingly. The resulting word contains only $\mathbf{1}$s (indicated by $z_{\mathbf{1}}$). Overall, this corresponds to $\mathbf{0}\mathbf{0}\mathbf{0}\mathbf{0}\Rightarrow_{\mathfrak{R}}\mathbf{0}%
\mathbf{2}\mathbf{2}\mathbf{0}\Rightarrow_{\mathfrak{R}}\mathbf{1}\mathbf{1}%
\mathbf{2}\mathbf{0}\Rightarrow_{\mathfrak{R}}\mathbf{1}\mathbf{1}\mathbf{1}%
\mathbf{1}$.\par
 
% END Thmtheorem30.p1
 
% END Thmtheorem30
 
% BEGIN S3.SSx2.SSS0.P0.SPx3.p2
 
Having only “well-formed” terms to consider (cf. Example 10 and Example 11), we can focus on the functional properties of rewriting.\par
 
% END S3.SSx2.SSS0.P0.SPx3.p2
 
% END S3.SSx2.SSS0.P0.SPx3
 
% END S3.SSx2
 
% BEGIN S3.SSx3
 
\par\medskip\noindent\textbf{Semantic Constraints}\par
 
% BEGIN S3.SSx3.p1
 
We formulate typed terms (Definition 31) and $\beta$-equivalence constraints characterizing word expansion (Lemma 36) and rewriting (Lemma 34). The presented terms are programs which realize the intended meaning (Example 30) of “well-formed” terms in $\mathcal{Q}_{m}$ and $\mathcal{R}_{m}$. Specifically, $G_{\star}$ realizes word expansion, $G_{\mathbf{0}}$ realizes initialization with $\mathbf{0}$s, $G_{ab\Rightarrow cd}$ realizes rule application, and $G_{j}^{i}$ controls the effect at position $i$ for rule application at position $j$. A word of length $m$ is represented by $m+1$ right-hand sides of $\beta$-equivalences.\par
 
% END S3.SSx3.p1
 
% BEGIN Thmtheorem31
 
\par\medskip\noindent\textbf{Definition 31 (Typed Terms $G_{\star}$, $G_{\mathbf{0}}$, $G_{ab\Rightarrow cd}$, $G_{i}^{j}$).}\par
 
% BEGIN Thmtheorem31.p1
    $\displaystyle G_{\star}:=$ $\displaystyle\lambda h.\lambda g.\lambda s_{\mathbf{0}}s_{\mathbf{1}}\ldots s_%
{\mathbf{K}}s_{\$}s_{\bullet}s_{\top}s_{\bot}.\textnormal{{case }}h\,\pi_{\top%
}\textnormal{{ of }}\langle s_{\bot}$      $\displaystyle\hskip 10.00002pt\mid\bullet\mapsto\textnormal{{case }}{g\,\delta%
_{\bullet}}\textnormal{{ of }}\langle s_{\bot}\mid\mathbf{0}\mapsto s_{\mathbf%
{0}}\mid\$\mapsto\textnormal{{case }}{g\,\delta_{\mathbf{0}}}\textnormal{{ of %
}}\langle s_{\bot}\mid\mathbf{1}\mapsto s_{\$}\rangle\rangle$      $\displaystyle\hskip 10.00002pt\mid\mathbf{0}\mapsto\textnormal{{case }}{g\,%
\delta_{\mathbf{1}}}\textnormal{{ of }}\langle s_{\bot}\mid\mathbf{0}\mapsto s%
_{\mathbf{1}}\rangle$      $\displaystyle\hskip 10.00002pt\mid\mathbf{1}\mapsto\textnormal{{case }}{g\,%
\delta_{\bullet}}\textnormal{{ of }}\langle s_{\bot}\mid\mathbf{0}\mapsto s_{%
\mathbf{0}}\rangle\rangle$     $\displaystyle G_{\mathbf{0}}:=$ $\displaystyle\lambda h.\lambda x.\lambda s_{\mathbf{0}}s_{\mathbf{1}}\ldots s_%
{\mathbf{K}}s_{\$}s_{\bullet}s_{\top}s_{\bot}.\textnormal{{case }}h\,\pi_{\top%
}\textnormal{{ of }}\langle s_{\bot}$      $\displaystyle\quad\mid\bullet\mapsto\textnormal{{case }}{x}\textnormal{{ of }}%
\langle s_{\bot}\mid\mathbf{0}\mapsto s_{\mathbf{0}}\mid\mathbf{1}\mapsto s_{%
\$}\rangle$      $\displaystyle\quad\mid\mathbf{0}\mapsto\textnormal{{case }}{x}\textnormal{{ of%
 }}\langle s_{\bot}\mid\mathbf{0}\mapsto s_{\mathbf{1}}\rangle$      $\displaystyle\quad\mid\mathbf{1}\mapsto\textnormal{{case }}{x}\textnormal{{ of%
 }}\langle s_{\bot}\mid\mathbf{0}\mapsto s_{\mathbf{0}}\rangle\rangle$     $\displaystyle G_{ab\Rightarrow cd}:=$ $\displaystyle\lambda h.\lambda x.\lambda s_{\mathbf{0}}s_{\mathbf{1}}\ldots s_%
{\mathbf{K}}s_{\$}s_{\bullet}s_{\top}s_{\bot}.\textnormal{{case }}h\,\pi_{\top%
}\textnormal{{ of }}\langle s_{\bot}$      $\displaystyle\quad\mid\bullet\mapsto x\,s_{\mathbf{0}}\,s_{\mathbf{1}}\,\ldots
s%
_{\mathbf{K}}\,s_{\$}\,s_{\bullet}\,s_{\top}\,s_{\bot}$      $\displaystyle\quad\mid\mathbf{0}\mapsto\textnormal{{case }}{x}\textnormal{{ of%
 }}\langle s_{\bot}\mid d\mapsto s_{b}\rangle$      $\displaystyle\quad\mid\mathbf{1}\mapsto\textnormal{{case }}{x}\textnormal{{ of%
 }}\langle s_{\bot}\mid c\mapsto s_{a}\rangle\rangle$     $\displaystyle G_{j}^{i}:=$ $\displaystyle\begin{cases}\delta_{\mathbf{1}}&\text{if }i=j\\
\delta_{\mathbf{0}}&\text{if }i=j+1\\
\delta_{\bullet}&\text{else}\end{cases}$       \[\emptyset\vdash G_{\star}:\Gamma_{m}(z_{\star})\qquad\quad\emptyset\vdash G_{%
\mathbf{0}}:\Gamma_{m}(z_{\mathbf{0}})\qquad\quad\emptyset\vdash G_{ab%
\Rightarrow cd}:\Gamma_{m}(r_{i})\text{ for }i\in\{1,\ldots,L\}\]    
% END Thmtheorem31.p1
 
% END Thmtheorem31
 
% BEGIN S3.SSx3.p2
 
Similarly to substitutions $S_{F}$ and $S_{H}$, we introduce the following substitution $S_{G}$.\par
 
% END S3.SSx3.p2
 
% BEGIN Thmtheorem32
 
\par\medskip\noindent\textbf{Definition 32 (Substitution $S_{G}$).}\par
 
% BEGIN Thmtheorem32.p1
 
\par
    \[S_{G}(z_{\star}):=G_{\star}\qquad S_{G}(z_{\mathbf{1}}):=\pi_{\mathbf{1}}%
\qquad S_{G}(z_{\mathbf{0}}):=G_{\mathbf{0}}\qquad S_{G}(r_{j}):=G_{R_{j}}%
\text{ for }j\in{1,\ldots L}\]    
% END Thmtheorem32.p1
 
% END Thmtheorem32
 
% BEGIN S3.SSx3.p3
 
The following Example 33 illustrates how the term $G_{ab\Rightarrow cd}$ represents rule application.\par
 
% END S3.SSx3.p3
 
% BEGIN Thmtheorem33
 
\par\medskip\noindent\textbf{Example 33.}\par
 
% BEGIN Thmtheorem33.p1
 
Consider for symbols $\mathbf{0},\mathbf{1},\ldots,\mathbf{5}$ an application of the rule $\mathbf{1}\mathbf{2}\Rightarrow\mathbf{4}\mathbf{5}$ at position $2$ in order to rewrite the word $\mathbf{0}\mathbf{1}\mathbf{2}\mathbf{3}$ to $\mathbf{0}\mathbf{4}\mathbf{5}\mathbf{3}$. Accordingly, we have:\par
  Position $1$:  
% BEGIN S3.I2.ix1.p1
 
$G_{\mathbf{1}\mathbf{2}\Rightarrow\mathbf{4}\mathbf{5}}\,G_{2}^{1}\,\pi_{%
\mathbf{0}}=_{\beta}\pi_{\mathbf{0}}$\par
 
% END S3.I2.ix1.p1
  Position $2$:  
% BEGIN S3.I2.ix2.p1
 
$G_{\mathbf{1}\mathbf{2}\Rightarrow\mathbf{4}\mathbf{5}}\,G_{2}^{2}\,\pi_{%
\mathbf{4}}=_{\beta}\pi_{\mathbf{1}}$\par
 
% END S3.I2.ix2.p1
  Position $3$:  
% BEGIN S3.I2.ix3.p1
 
$G_{\mathbf{1}\mathbf{2}\Rightarrow\mathbf{4}\mathbf{5}}\,G_{2}^{3}\,\pi_{%
\mathbf{5}}=_{\beta}\pi_{\mathbf{2}}$\par
 
% END S3.I2.ix3.p1
  Position $4$:  
% BEGIN S3.I2.ix4.p1
 
$G_{\mathbf{1}\mathbf{2}\Rightarrow\mathbf{4}\mathbf{5}}\,G_{2}^{4}\,\pi_{%
\mathbf{3}}=_{\beta}\pi_{\mathbf{3}}$\par
 
% END S3.I2.ix4.p1
   
% END Thmtheorem33.p1
 
% END Thmtheorem33
 
% BEGIN S3.SSx3.p4
 
The above observation is generalized for terms in $\mathcal{Q}_{m}$ in the following Lemma 34.\par
 
% END S3.SSx3.p4
 
% BEGIN Thmtheorem34
 
\par\medskip\noindent\textbf{Lemma 34.}\par
 
% BEGIN Thmtheorem34.p1
 
For $m>0$, let $a_{1},\ldots,a_{m+1}\in\{\mathbf{0},\mathbf{1},\ldots,\mathbf{K}\}$ and $M\in\mathcal{Q}_{m}$. If $\Gamma_{m}\vdash M:\kappa$, $S_{G}(M)[p_{1}:=G_{1}^{0},\ldots,p_{m}:=G_{m}^{0}]=_{\beta}\pi_{\mathbf{1}}$, and $S_{G}(M)[p_{1}:=G_{1}^{i},\ldots,p_{m}:=G_{m}^{i}]=_{\beta}\pi_{a_{i}}$ for $i\in\{1,\ldots,m+1\}$, then $a_{1}\ldots a_{m+1}\Rightarrow_{\mathfrak{R}}^{*}\mathbf{1}^{m+1}$.\par
 
% END Thmtheorem34.p1
 
% END Thmtheorem34
 
% BEGIN S3.SSx3.SSS0.P0.SPx1
 
\par\medskip\noindent\textbf{Proof.}\par
 
% BEGIN S3.SSx3.SSS0.P0.SPx1.p1
 
Induction on the size of $M$ and case analysis using Definition 22. $\hfill\blacktriangleleft$\par
 
% END S3.SSx3.SSS0.P0.SPx1.p1
 
% BEGIN S3.SSx3.SSS0.P0.SPx1.p2
 
The following Example 35 builds upon the previous Example 30 and illustrates the intended meaning (rewriting $\mathbf{0}$s to $\mathbf{1}$s) of a “well-formed” example term in $\mathcal{Q}_{3}$.\par
 
% END S3.SSx3.SSS0.P0.SPx1.p2
 
% BEGIN Thmtheorem35
 
\par\medskip\noindent\textbf{Example 35.}\par
 
% BEGIN Thmtheorem35.p1
 
Assume $\mathfrak{R}=\{\mathbf{0}\mathbf{0}\Rightarrow\mathbf{2}\mathbf{2},\mathbf{0}%
\mathbf{2}\Rightarrow\mathbf{1}\mathbf{1},\mathbf{2}\mathbf{0}\Rightarrow%
\mathbf{1}\mathbf{1}\}$ over the alphabet $\{\mathbf{0},\mathbf{1},\mathbf{2}\}$, and consider the term $N=r_{1}\,p_{2}\,(r_{2}\,p_{1}\,(r_{3}\,p_{3}\,z_{\mathbf{1}}))$ from Example 30. Replacing $G_{j}^{i}$ accordingly for $i\in\{0,\ldots,5\}$ and $j\in\{1,2,3\}$, we have the following $\beta$-equivalences $(0)$ – $(4)$.\par
    \[\begin{array}[]{lll}(0)&&S_{G}(N)[p_{1}:=\delta_{\bullet},p_{2}:=\delta_{%
\bullet},p_{3}:=\delta_{\bullet}]=_{\beta}\pi_{\mathbf{1}}\\
(1)&&S_{G}(N)[p_{1}:=\delta_{\mathbf{1}},p_{2}:=\delta_{\bullet},p_{3}:=\delta%
_{\bullet}]=_{\beta}\pi_{\mathbf{0}}\\
(2)&&S_{G}(N)[p_{1}:=\delta_{\mathbf{0}},p_{2}:=\delta_{\mathbf{1}},p_{3}:=%
\delta_{\bullet}]=_{\beta}\pi_{\mathbf{0}}\\
(3)&&S_{G}(N)[p_{1}:=\delta_{\bullet},p_{2}:=\delta_{\mathbf{0}},p_{3}:=\delta%
_{\mathbf{1}}]=_{\beta}\pi_{\mathbf{0}}\\
(4)&&S_{G}(N)[p_{1}:=\delta_{\bullet},p_{2}:=\delta_{\bullet},p_{3}:=\delta_{%
\mathbf{0}}]=_{\beta}\pi_{\mathbf{0}}\end{array}\]    
% END Thmtheorem35.p1
  
% BEGIN Thmtheorem35.p2
 
In accordance with Lemma 34, we have that $\mathbf{0}^{4}\Rightarrow_{\mathfrak{R}}^{*}\mathbf{1}^{4}$. Equivalences $(1)$ – $(4)$ witness the particular rewriting steps at positions $1$ – $4$ (cf. Example 30 and Example 33).\par
 
% END Thmtheorem35.p2
 
% END Thmtheorem35
 
% BEGIN S3.SSx3.SSS0.P0.SPx1.p3
 
Complementarily to word rewriting, the following Lemma 36 characterizes word expansion and initialization with $\mathbf{0}$s.\par
 
% END S3.SSx3.SSS0.P0.SPx1.p3
 
% BEGIN Thmtheorem36
 
\par\medskip\noindent\textbf{Lemma 36.}\par
 
% BEGIN Thmtheorem36.p1
 
If $M\in\mathcal{R}_{1}$ such that $\Gamma_{1}\vdash M:\kappa$, $S_{G}(M)[p_{1}:=\delta_{\bullet}]=_{\beta}\pi_{\$}$, $S_{G}(M)[p_{1}:=\delta_{\mathbf{1}}]=_{\beta}\pi_{\mathbf{0}}$, and $S_{G}(M)[p_{1}:=\delta_{\mathbf{0}}]=_{\beta}\pi_{\mathbf{1}}$, then there exists an $m>0$ and an $N\in\mathcal{Q}_{m}$ such that $\Gamma_{m}\vdash N:\kappa$, $S_{G}(N)[p_{1}:=G_{1}^{0},\ldots,p_{m}:=G_{m}^{0}]=_{\beta}\pi_{\mathbf{1}}$, and $S_{G}(N)[p_{1}:=G_{1}^{i},\ldots,p_{m}:=G_{m}^{i}]=_{\beta}\pi_{\mathbf{0}}$ for $i\in\{1,\ldots,m+1\}$.\par
 
% END Thmtheorem36.p1
 
% END Thmtheorem36
 
% END S3.SSx3.SSS0.P0.SPx1
 
% BEGIN S3.SSx3.SSS0.P0.SPx2
 
\par\medskip\noindent\textbf{Proof.}\par
 
% BEGIN S3.SSx3.SSS0.P0.SPx2.p1
 
Considering the general case $M\in\mathcal{R}_{m^{\prime}}$ for $m^{\prime}>0$, induction on the size of $M$ and case analysis using Definition 22. $\hfill\blacktriangleleft$\par
 
% END S3.SSx3.SSS0.P0.SPx2.p1
 
% BEGIN S3.SSx3.SSS0.P0.SPx2.p2
 
The following Example 37 complements the previous Example 35 and illustrates the intended meaning (word expansion and initialization with $\mathbf{0}$s) of a “well-formed” term in $\mathcal{R}_{1}$.\par
 
% END S3.SSx3.SSS0.P0.SPx2.p2
 
% BEGIN Thmtheorem37
 
\par\medskip\noindent\textbf{Example 37.}\par
 
% BEGIN Thmtheorem37.p1
 
Assume $\mathfrak{R}=\{\mathbf{0}\mathbf{0}\Rightarrow\mathbf{2}\mathbf{2},\mathbf{0}%
\mathbf{2}\Rightarrow\mathbf{1}\mathbf{1},\mathbf{2}\mathbf{0}\Rightarrow%
\mathbf{1}\mathbf{1}\}$ over the alphabet $\{\mathbf{0},\mathbf{1},\mathbf{2}\}$, and consider the terms $M_{1}:=z_{\star}\,p_{1}\,(\lambda p_{2}.M_{2})$, $M_{2}:=z_{\star}\,p_{2}\,(\lambda p_{3}.M_{3})$, $M_{3}:=z_{\mathbf{0}}\,p_{3}\,N$, and $N:=r_{1}\,p_{2}\,(r_{2}\,p_{1}\,(r_{3}\,p_{3}\,z_{\mathbf{1}}))$ from Example 35. Proceeding bottom up, we have $M_{3}\in\mathcal{R}_{3}$ and the following $\beta$-equivalences hold:\par
    \[\begin{array}[]{l}S_{G}(M_{3})[p_{1}:=\delta_{\bullet},p_{2}:=\delta_{\bullet}%
,p_{3}:=\delta_{\bullet}]=_{\beta}\pi_{\$}\\
S_{G}(M_{3})[p_{1}:=\delta_{\mathbf{1}},p_{2}:=\delta_{\bullet},p_{3}:=\delta_%
{\bullet}]=_{\beta}\pi_{\mathbf{0}}\\
S_{G}(M_{3})[p_{1}:=\delta_{\mathbf{0}},p_{2}:=\delta_{\mathbf{1}},p_{3}:=%
\delta_{\bullet}]=_{\beta}\pi_{\mathbf{0}}\\
S_{G}(M_{3})[p_{1}:=\delta_{\bullet},p_{2}:=\delta_{\mathbf{0}},p_{3}:=\delta_%
{\mathbf{1}}]=_{\beta}\pi_{\mathbf{0}}\\
S_{G}(M_{3})[p_{1}:=\delta_{\bullet},p_{2}:=\delta_{\bullet},p_{3}:=\delta_{%
\mathbf{0}}]=_{\beta}\pi_{\mathbf{1}}\end{array}\]    
Additionally, $M_{2}\in\mathcal{R}_{2}$, $M_{1}\in\mathcal{R}_{1}$, and the following $\beta$-equivalences hold:\par
    \[\begin{array}[]{l}S_{G}(M_{2})[p_{1}:=\delta_{\bullet},p_{2}:=\delta_{\bullet}%
]=_{\beta}\pi_{\$}\qquad\qquad S_{G}(M_{1})[p_{1}:=\delta_{\bullet}]=_{\beta}%
\pi_{\$}\\
S_{G}(M_{2})[p_{1}:=\delta_{\mathbf{1}},p_{2}:=\delta_{\bullet}]=_{\beta}\pi_{%
\mathbf{0}}\qquad\qquad S_{G}(M_{1})[p_{1}:=\delta_{\mathbf{1}}]=_{\beta}\pi_{%
\mathbf{0}}\\
S_{G}(M_{2})[p_{1}:=\delta_{\mathbf{0}},p_{2}:=\delta_{\mathbf{1}}]=_{\beta}%
\pi_{\mathbf{0}}\qquad\qquad S_{G}(M_{1})[p_{1}:=\delta_{\mathbf{0}}]=_{\beta}%
\pi_{\mathbf{1}}\\
S_{G}(M_{2})[p_{1}:=\delta_{\bullet},p_{2}:=\delta_{\mathbf{0}}]=_{\beta}\pi_{%
\mathbf{1}}\end{array}\]    
% END Thmtheorem37.p1
 
% BEGIN Thmtheorem37.p2
 
In combination with the previous Example 35, the term $M_{1}\in\mathcal{R}_{1}$ represents word expansion up to length $4$, followed by initialization with $\mathbf{0}$s, and rewriting to $\mathbf{1}$s.\par
 
% END Thmtheorem37.p2
 
% END Thmtheorem37
 
% BEGIN S3.SSx3.SSS0.P0.SPx2.p3
 
Next, we combine syntactic and semantic constraints in the following key Lemma 38.\par
 
% END S3.SSx3.SSS0.P0.SPx2.p3
 
% BEGIN Thmtheorem38
 
\par\medskip\noindent\textbf{Lemma 38.}\par
 
% BEGIN Thmtheorem38.p1
 
There exists an $n\in\mathbb{N}$ such that $\mathbf{0}^{n+1}\Rightarrow_{\mathfrak{R}}^{*}\mathbf{1}^{n+1}$ iff there exists a term $M$ in normal form, such that the following conditions hold:\par
 
% END Thmtheorem38.p1
 
% BEGIN Thmtheorem38.p2
 
\begin{center}\includegraphics[width=0.95\linewidth,height=0.45\textheight,keepaspectratio]{formula-03.pdf}\end{center}
 
% END Thmtheorem38.p2
 
% BEGIN Thmtheorem38.p3
 
\begin{center}\includegraphics[width=0.95\linewidth,height=0.45\textheight,keepaspectratio]{formula-04.pdf}\end{center}
 
% END Thmtheorem38.p3
 
% END Thmtheorem38
 
% END S3.SSx3.SSS0.P0.SPx2
 
% BEGIN S3.SSx3.SSS0.P0.SPx3
 
\par\medskip\noindent\textbf{Proof.}\par
 
% BEGIN S3.SSx3.SSS0.P0.SPx3.p1
 
The direction from left to right proceeds in two steps. First, by induction on the number of rewriting steps we construct a term $N\in\mathcal{Q}_{n}$ (easy converse of Lemma 34). Second, by induction on $n$ we construct a term $M^{\prime}\in\mathcal{R}_{1}$ containing $N\in\mathcal{Q}_{n}$ as a subterm (easy converse of Lemma 36). Then, the solution is $M:=\lambda r_{1}\ldots r_{L}.\lambda z_{\mathbf{0}}z_{\mathbf{1}}z_{\star}p_{1%
}.M^{\prime}$.\par
 
% END S3.SSx3.SSS0.P0.SPx3.p1
 
% BEGIN S3.SSx3.SSS0.P0.SPx3.p2
 
The direction from right to left proceeds in two steps. First, by Theorem 29 we have $\lambda r_{1}\ldots r_{L}.\lambda z_{\mathbf{0}}z_{\mathbf{1}}z_{\star}p_{1}.M%
^{\prime}$ for some $M^{\prime}\in\mathcal{R}_{1}$. Second, by Lemma 36 and Lemma 34 we have $\mathbf{0}^{n+1}\Rightarrow_{\mathfrak{R}}^{*}\mathbf{1}^{n+1}$ for some $n\in\mathbb{N}$. $\hfill\blacktriangleleft$\par
 
% END S3.SSx3.SSS0.P0.SPx3.p2
 
% BEGIN S3.SSx3.SSS0.P0.SPx3.p3
 
Finally, we present the combination of constraints from the above Lemma 38 as a matching instance $F_{\mathfrak{R}}\,\mathsf{X}=_{\beta}N_{\mathfrak{R}}$. This constitutes the main result of the present work.\par
 
% END S3.SSx3.SSS0.P0.SPx3.p3
 
% BEGIN Thmtheorem39
 
\par\medskip\noindent\textbf{Theorem 39.}\par
 
% BEGIN Thmtheorem39.p1
 
Problem $\mathbf{0}^{+}\Rightarrow^{*}\mathbf{1}^{+}$ many-one reduces to higher-order $\beta$-matching.\par
 
% END Thmtheorem39.p1
 
% END Thmtheorem39
 
% END S3.SSx3.SSS0.P0.SPx3
 
% BEGIN S3.SSx3.SSS0.P0.SPx4
 
\par\medskip\noindent\textbf{Proof.}\par
 
% BEGIN S3.SSx3.SSS0.P0.SPx4.p1
 
Given a simple semi-Thue system $\mathfrak{R}=\{R_{1},\ldots,R_{L}\}$ due to Lemma 38 there exists an $n\in\mathbb{N}$ such that $\mathbf{0}^{n+1}\Rightarrow_{\mathfrak{R}}^{*}\mathbf{1}^{n+1}$ iff the instance $F_{\mathfrak{R}}\,\mathsf{X}=_{\beta}N_{\mathfrak{R}}$ of higher-order $\beta$-matching is solvable, where\par
  
\par\noindent  $\blacksquare$ 
% BEGIN S3.I3.i1.p1
 
\begin{center}\includegraphics[width=0.95\linewidth,height=0.45\textheight,keepaspectratio]{formula-05.pdf}\end{center}
 
% END S3.I3.i1.p1
 
 
\par\noindent  $\blacksquare$ 
% BEGIN S3.I3.i2.p1
 
$N_{\mathfrak{R}}:=\lambda y.y\,(\lambda u.u)\,\pi_{\$}\,\pi_{\$}\,\pi_{\mathbf%
{0}}\,\pi_{\mathbf{1}}$\par
 
% END S3.I3.i2.p1
 
 
\par\noindent  $\blacksquare$ 
% BEGIN S3.I3.i3.p1
 
$\sigma_{\mathfrak{R}}:=\Gamma_{1}(r_{1})\to\cdots\to\Gamma_{1}(r_{L})\to\Gamma%
_{1}(z_{\mathbf{0}})\to\Gamma_{1}(z_{\mathbf{1}})\to\Gamma_{1}(z_{\star})\to%
\Gamma_{1}(p_{1})\to\kappa$\par
 
% END S3.I3.i3.p1
 
 
\par\noindent  $\blacksquare$ 
% BEGIN S3.I3.i4.p1
 
$\tau_{\mathfrak{R}}:=((\kappa\to\kappa)\to\kappa\to\kappa\to\kappa\to\kappa\to%
\iota)\to\iota$\par
 
% END S3.I3.i4.p1
 
 
\par\noindent  $\blacksquare$ 
% BEGIN S3.I3.i5.p1
 
$\emptyset\vdash F_{\mathfrak{R}}:\sigma_{\mathfrak{R}}\to\tau_{\mathfrak{R}}$\par
 
% END S3.I3.i5.p1
 
 
\par\noindent  $\blacksquare$ 
% BEGIN S3.I3.i6.p1
 
$\emptyset\vdash N_{\mathfrak{R}}:\tau_{\mathfrak{R}}$\par
 
% END S3.I3.i6.p1
 
  
$\hfill\blacktriangleleft$\par
 
% END S3.SSx3.SSS0.P0.SPx4.p1
 
% BEGIN Thmtheorem40
 
\par\medskip\noindent\textbf{Theorem 40.}\par
 
% BEGIN Thmtheorem40.p1
 
Higher-order $\beta$-matching (Problem 7) is undecidable.\par
 
% END Thmtheorem40.p1
 
% END Thmtheorem40
 
% END S3.SSx3.SSS0.P0.SPx4
 
% BEGIN S3.SSx3.SSS0.P0.SPx5
 
\par\medskip\noindent\textbf{Proof.}\par
 
% BEGIN S3.SSx3.SSS0.P0.SPx5.p1
 
By reduction from the undecidable Problem $\mathbf{0}^{+}\Rightarrow^{*}\mathbf{1}^{+}$ (Theorem 17 and Theorem 39). $\hfill\blacktriangleleft$\par
 
% END S3.SSx3.SSS0.P0.SPx5.p1
 
% END S3.SSx3.SSS0.P0.SPx5
 
% END S3.SSx3
 
% END S3


% BEGIN S4
 
\par\medskip\noindent\textbf{4 Mechanization}\par
 
% BEGIN S4.p1
 
This section provides a brief overview over the mechanization of undecidability of higher-order $\beta$-matching (Theorem 40) using the Coq proof assistant [19]. The mechanization is constructive (in the technical sense of axiom-free Coq) and spans approximately $4000$ lines of code, consisting of the following parts:\par
  
\par\noindent  $\blacksquare$ 
% BEGIN S4.I1.i1.p1
 
HOMatching.v contains definitions of the simply typed $\lambda$-calculus [ ] and higher-order $\beta$-matching [ ].\par
 
% END S4.I1.i1.p1
 
 
\par\noindent  $\blacksquare$ 
% BEGIN S4.I1.i2.p1
 
Util/stlc\_facts.v and Util/term\_facts.v contain basic properties of the simply typed $\lambda$-calculus, such as confluence of $\beta$-reduction [ ], substitution lemmas [ ], and type preservation properties [ ].\par
 
% END S4.I1.i2.p1
 
 
\par\noindent  $\blacksquare$ 
% BEGIN S4.I1.i3.p1
 
Reductions/SSTS01\_to\_HOMbeta.v contains the reduction from Problem $\mathbf{0}^{+}\Rightarrow^{*}\mathbf{1}^{+}$ to higher-order $\beta$-matching [ ].\par
 
% END S4.I1.i3.p1
 
 
\par\noindent  $\blacksquare$ 
% BEGIN S4.I1.i4.p1
 
HOMatching\_undec.v contains the undecidability result for higher-order $\beta$-matching [ ].\par
 
% END S4.I1.i4.p1
 
  
% END S4.p1
 
% BEGIN S4.p2
 
The simple type system stlc is mechanized in HOMatching.v, borrowing the existing term definitions from the library [ ], for which variable binding is addressed via the unscoped de Bruijn approach [4]. The proposition stlc Gamma M t mechanizes that the term M is assigned the simple type t in the simple type environment Gamma.\par
  \ensuremath{\downarrow} 
% BEGIN lstnumberx1
 Inductive ty : Type := 
% END lstnumberx1
 
% BEGIN lstnumberx2
  | atom (* type variable *) 
% END lstnumberx2
 
% BEGIN lstnumberx3
  | arr (s t : ty). (* function type *) 
% END lstnumberx3
 
% BEGIN lstnumberx4
 
% END lstnumberx4
 
% BEGIN lstnumberx5
 Inductive term : Type := 
% END lstnumberx5
 
% BEGIN lstnumberx6
  | var (n : nat) : term (* term variable *) 
% END lstnumberx6
 
% BEGIN lstnumberx7
  | app (s : term) (t : term) : term (* application *) 
% END lstnumberx7
 
% BEGIN lstnumberx8
  | lam (s : term). (* abstraction *) 
% END lstnumberx8
 
% BEGIN lstnumberx9
 
% END lstnumberx9
 
% BEGIN lstnumberx10
 Inductive stlc (Gamma : list ty) : term -> ty -> Prop := 
% END lstnumberx10
 
% BEGIN lstnumberx11
  | stlc\_var x t : nth\_error Gamma x = Some t -> 
% END lstnumberx11
 
% BEGIN lstnumberx12
  stlc Gamma (var x) t (* variable rule *) 
% END lstnumberx12
 
% BEGIN lstnumberx13
  | stlc\_app M N s t : stlc Gamma M (arr s t) -> stlc Gamma N s -> 
% END lstnumberx13
 
% BEGIN lstnumberx14
  stlc Gamma (app M N) t (* application rule *) 
% END lstnumberx14
 
% BEGIN lstnumberx15
  | stlc\_lam M s t : stlc (cons s Gamma) M t -> 
% END lstnumberx15
 
% BEGIN lstnumberx16
  stlc Gamma (lam M) (arr s t). (* abstraction rule *) 
% END lstnumberx16
  
% END S4.p2
 
% BEGIN S4.p3
 
Higher-order $\beta$-matching is mechanized as the predicate HOMbeta: given terms F of type arr s t and N of type t, is there a simply typed term M of type s such that app F M is $\beta$-equivalent (reflexive, symmetric, transitive closure of step) to N?\par
  \ensuremath{\downarrow} 
% BEGIN lstnumberx17
 Definition HOMbeta : \{ ’(s, t, F, N) : (ty * ty * term * term) 
% END lstnumberx17
 
% BEGIN lstnumberx18
  | stlc nil F (arr s t) /\textbackslash{} stlc nil N t \} -> Prop := 
% END lstnumberx18
 
% BEGIN lstnumberx19
  fun ’(exist \_ (s, t, F, N) \_) => 
% END lstnumberx19
 
% BEGIN lstnumberx20
  exists (M : term), stlc nil M s /\textbackslash{} 
% END lstnumberx20
 
% BEGIN lstnumberx21
  clos\_refl\_sym\_trans term step (app F M) N. 
% END lstnumberx21
  
% END S4.p3
 
% BEGIN S4.p4
 
The proposition undecidable HOMbeta [ ] mechanizes the undecidability of the predicate HOMbeta, relying on the following library definition [7, Chapter 19]. A predicate p is undecidable, if existence of a computable decider for p implies recursive co-enumerability of the (Turing machine) Halting Problem.\par
  \ensuremath{\downarrow} 
% BEGIN lstnumberx22
 Definition undecidable \{X\} (p : X -> Prop) := 
% END lstnumberx22
 
% BEGIN lstnumberx23
  decidable p -> enumerable (complement SBTM\_HALT). 
% END lstnumberx23
  
Since the Halting Problem is recursively enumerable, decidability of p would imply decidability of the Halting Problem.\par
 
% END S4.p4
 
% END S4
